
\subsubsection{4.1 Research Questions}

\textbf{RQ1} (h-e1): Does aggregation function choice produce statistically different AUROC values?  
\textbf{RQ2} (h-m1): Do recall-failure hallucinations produce peaked token probability distributions?  
\textbf{RQ3} (h-m2): Does Spearman $\rho$ confirm directional advantage across all four (model $\times$ dataset) pairs?  
\textbf{RQ4} (h-m3): Do AUROC differentials confirm P1 and P2 with $\geq 0.02$ margin and bootstrap CI excluding zero?

\subsubsection{4.2 Datasets}

\begin{table*}[htbp]
\centering
\resizebox{\dimexpr 2\columnwidth+\columnsep\relax}{!}{%
\begin{tabular}{lllll}
\toprule
Dataset & Hallucination Type & N (LLaMA / Mistral) & Label Protocol & Source \\
\midrule
TriviaQA & Recall-failure & 488 / 476 & Exact match & Farquhar 2024 splits (jlko/semantic\_uncertainty) \\
TruthfulQA (gen.) & Imitative-falsehood & 810 / 796 & ROUGE-L $\geq$ 0.3 & HuggingFace generation subset \\
\bottomrule
\end{tabular}
}
\end{table*}

\textbf{TriviaQA} \citep{Joshi2017TriviaQA} uses the Farquhar 2024 \texttt{jlko/semantic\textbackslash \_uncertainty} splits, enabling direct comparison to published baselines. Minor sample count variation between models reflects tokenization differences in empty-generation filtering. \textbf{TruthfulQA} \citep{Lin2021TruthfulQA} uses the HuggingFace generation subset; binary labels are assigned at inference time via ROUGE-L $\geq$ 0.3 against \texttt{best\textbackslash \_answer} and \texttt{correct\textbackslash \_answers} fields, following Fadeeva et al. [2024]. \textbf{Natural Questions (NQ)} was planned but could not be evaluated: inference ran during h-e1 but the score cache file was not persisted to disk; this is a data availability gap, not a mechanism failure.

\subsubsection{4.3 Implementation Details}

Inference: float16 precision, \texttt{flash\textbackslash \_attention\textbackslash \_2} with fallback, \texttt{device\textbackslash \_map=auto}, single NVIDIA H100 NVL GPU (24 GB VRAM), batch size 1. Bootstrap CI: n=1000 resamples, seed=42, percentile method, paired=True (\texttt{scipy.stats.bootstrap}). Statistical tests for h-m1: two-sample t-test (\texttt{scipy.stats.ttest\textbackslash \_ind}). Peakedness metric: \texttt{max(|token\textbackslash \_logprobs|) / mean(|token\textbackslash \_logprobs|)}.

\subsubsection{4.4 Literature Baselines (Reference, Cross-Pipeline)}

\begin{table}[htbp]
\centering
\resizebox{\columnwidth}{!}{%
\begin{tabular}{lll}
\toprule
Method & Inference Cost & Published AUROC (TriviaQA) \\
\midrule
Mean log-prob (CCP) [Fadeeva et al., 2024] & 1 forward pass & 0.72--0.80 \\
Predictive entropy (sum) [Farquhar et al., 2024] & 1 forward pass & $\approx 0.72$ \\
Semantic Entropy [Farquhar et al., 2024] & 10 forward passes & $\approx 0.79$ \\
SelfCheckGPT (NLI) [Manakul et al., 2023] & 20 forward passes & 0.72--0.78 \\
\bottomrule
\end{tabular}
}
\end{table}

Direct numerical comparison to these baselines is subject to cross-pipeline caveats (different hardware, dataset splits, tokenization, and evaluation protocols); the magnitude of observed differences is discussed in Section 6.3.

